\documentclass[10pt,letterpaper]{article}
\usepackage[margin=1in]{geometry}

\usepackage[utf8]{inputenc}

\usepackage{amsmath,amssymb}
\usepackage[expansion=false]{microtype}
\usepackage{booktabs,longtable,array,calc}
\usepackage[small,compact]{titlesec}
\usepackage{enumitem}
\usepackage{xcolor}
\usepackage{seqsplit}
\usepackage{needspace}
\usepackage[colorlinks=true,linkcolor=black,citecolor=black,urlcolor=blue!45!black]{hyperref}
\usepackage{fancyhdr}
\DeclareUnicodeCharacter{03C6}{\ensuremath{\varphi}}
\DeclareUnicodeCharacter{2212}{\ensuremath{-}}
\DeclareUnicodeCharacter{0394}{\ensuremath{\Delta}}
\DeclareUnicodeCharacter{2265}{\ensuremath{\geq}}
\DeclareUnicodeCharacter{2264}{\ensuremath{\leq}}
\DeclareUnicodeCharacter{00D7}{\ensuremath{\times}}
\DeclareUnicodeCharacter{00B1}{\ensuremath{\pm}}
\DeclareUnicodeCharacter{2208}{\ensuremath{\in}}
\DeclareUnicodeCharacter{03B1}{\ensuremath{\alpha}}
\DeclareUnicodeCharacter{2192}{\ensuremath{\rightarrow}}
\DeclareUnicodeCharacter{00A7}{\ensuremath{\mathsection}}
\DeclareUnicodeCharacter{2026}{\ldots}
\DeclareUnicodeCharacter{00B7}{\ensuremath{\cdot}}
\DeclareUnicodeCharacter{03C1}{\ensuremath{\rho}}
\newcommand{\usd}{{\usefont{OT1}{cmr}{m}{n}\char36}}
\renewcommand{\labelitemi}{$\bullet$}
\renewcommand{\labelitemii}{$\circ$}
\providecommand{\tightlist}{\setlength{\itemsep}{0pt}\setlength{\parskip}{0pt}}
\setlist{nosep,leftmargin=1.4em,topsep=2pt}
\setlength{\parskip}{0.38em}
\setlength{\parindent}{0pt}
\setcounter{secnumdepth}{0}
\setlength{\emergencystretch}{1.5em}
\titlespacing*{\section}{0pt}{1.1em}{0.35em}
\titleformat{\section}{\large\bfseries}{}{0pt}{}
\setlength{\LTpre}{6pt}\setlength{\LTpost}{4pt}
\renewcommand{\arraystretch}{1.14}
\setlength{\tabcolsep}{4.5pt}
\pagestyle{fancy}\fancyhf{}
\renewcommand{\headrulewidth}{0pt}
\fancyfoot[L]{\footnotesize\color{black!60}Seen and Unseen: Preregistration, v3.4}
\fancyfoot[R]{\footnotesize\color{black!60}\thepage}
\begin{document}
\begin{center}
{\Large\bfseries Seen and Unseen: Training Failure Share and\\[2pt] Trajectory Reuse in World-Model Evaluation}\\[6pt]
{\normalsize Preregistration \quad$\cdot$\quad v3.4 \quad$\cdot$\quad revised November 14, 2026}
\end{center}
\vspace{2pt}
This document preregisters three computational experiments.
It contains every confirmatory commitment: claims,
estimators, thresholds, exclusions and permitted fallbacks.

\needspace{4\baselineskip}
\hypertarget{summary}{%
\section{Summary}\label{summary}}

\textbf{Motivation.} Learned world models are increasingly
used to evaluate robot policies on episodes the model
generates (``in imagination'') instead of on real robots. If
such evaluations get the order of policies right but their
success rates wrong, teams can pick the right policy while
reporting the wrong number, and published comparisons can
mislead. Two properties of the data could cause such errors:
how many failures the world model saw in training, and
evaluation items that reuse training trajectories. On logged
Kuka data, a world model trained without failures estimated
about 51\% success against 36\% recorded (§G).

\textbf{Terms.} A world model is a learned simulator that
predicts the next observation from recent observations and
an action. In closed-loop evaluation, a policy acts on these
predictions for a whole episode (an imagined rollout), and a
judge (a learned classifier in the primary analysis) labels
the episode a success or failure. A real rollout runs the
same policy in the simulator from the same start. True
success is the simulator's own outcome for an episode. A
policy's true success is its rate over 10,000 simulator
episodes, and far pairs are pairs of policies at least 15
points apart in it. The headline is P1 and P3 together
(Experiment 1's row in the table below). A seed is one
training run. In Experiment 3, a block is one independently
trained forecaster with its matched continuations; it is
screened in if its base model beats constant-velocity
extrapolation. Errors are in arena widths, and the positive
control trains on the exact evaluation renders.

\textbf{Claims under test.} First, estimated success from
world-model policy evaluation depends on the share of
failures in the world model's training data, while policies
far apart in true success are misordered little more often
in imagined rollouts than by the same judge on real
rollouts. Second, in controlled 2D forecasting, reusing
training trajectories in evaluation lowers forecast error on
those trajectories relative to independent events.
Experiment 1 tests the first claim in closed loop.
Experiment 2 checks its failure-share part on logged robot
data. Experiment 3 tests the second claim.

\begingroup\small
\begin{longtable}[]{@{}
  >{\raggedright\arraybackslash}p{(\columnwidth - 8\tabcolsep) * \real{0.1600}}
  >{\raggedright\arraybackslash}p{(\columnwidth - 8\tabcolsep) * \real{0.2400}}
  >{\raggedright\arraybackslash}p{(\columnwidth - 8\tabcolsep) * \real{0.2000}}
  >{\raggedright\arraybackslash}p{(\columnwidth - 8\tabcolsep) * \real{0.2000}}
  >{\raggedright\arraybackslash}p{(\columnwidth - 8\tabcolsep) * \real{0.2000}}@{}}
\toprule\noalign{}
\begin{minipage}[b]{\linewidth}\raggedright
Experiment
\end{minipage} & \begin{minipage}[b]{\linewidth}\raggedright
Primary hypothesis
\end{minipage} & \begin{minipage}[b]{\linewidth}\raggedright
Test and margin
\end{minipage} & \begin{minipage}[b]{\linewidth}\raggedright
Supported means
\end{minipage} & \begin{minipage}[b]{\linewidth}\raggedright
Absent means
\end{minipage} \\
\midrule\noalign{}
\endhead
\bottomrule\noalign{}
\endlastfoot
1. Closed-loop evaluation in simulation (primary) & Raising
the training failure share from 10\% to 65\% lowers
estimated success, relative to the same judge on real
rollouts, by more than 3 points (P1); at a 10\% failure
share, the excess misordering of pairs at least 15 points
apart in true success (imagined minus real-rollout
misordering rate) is below 5 points (P3) & One-sided 95\%
bounds over 12 seeds; shift \ensuremath{>} 3 points; excess
misordering \ensuremath{<} 5 points & For these 12 policies,
estimates depend on the training failure share, and far
pairs are misordered less than 5 points more often than on
real rollouts & Failure share does not move closed-loop
estimates beyond ±3 points in this setting \\
2. Transition-matched control on logged Kuka data
(robustness check of a known drop; §G) & With every arm
trained on exactly 31,600 distinct transitions, raising the
failure share from 10\% to 65\% still lowers estimated
success (\(\Delta_A\)) & One-sided 95\% bound over 10 seeds;
drop \ensuremath{>} 3 points & The previously observed
open-loop drop (logged actions replayed in the world model)
survives when the number of training transitions is held
fixed & The drop is within ±3 points once the number of
transitions is held fixed \\
3. Trajectory reuse in controlled 2D forecasting & Reusing
training trajectories in evaluation lowers forecast error on
them relative to independent events (reuse gap G) &
Two-sided 95\% interval over the screened 16-channel blocks
(800 scheduled, at least 720 required); consequential if the
lower bound exceeds 0.010 arena widths; negligible if the
90\% interval lies inside ±0.002 and the positive control's
lower bound exceeds 0.002 & The reuse gap exceeds 0.010
arena widths, with lower error on the reused trajectories
themselves (source gain S \ensuremath{>} 0) and error on
independent events raised by less than 0.001 (harm H) & The
reuse gap is negligible, and training on the exact
evaluation renders (the positive control) lowers error on
them by more than 0.002 \\
\end{longtable}
\endgroup

P2 tests whether, at a 10\% failure share, imagined
estimates differ from the same judge's real-rollout
estimates by more than 5 points on average over policies
(mean absolute world-model shift). It is decided separately
from the headline, and a supported headline does not by
itself establish that this shift is large. P3 measures
misordering in excess of the same judge's on real rollouts:
if that judge itself misorders far pairs, a supported P3
shows limited additional ordering error in imagined
rollouts, not correct ranking (§D). Sections D--F give every
secondary hypothesis, fallback and failure rule.

\textbf{Margins and rationale.}

\begingroup\small
\begin{longtable}[]{@{}
  >{\raggedright\arraybackslash}p{(\columnwidth - 4\tabcolsep) * \real{0.2692}}
  >{\raggedright\arraybackslash}p{(\columnwidth - 4\tabcolsep) * \real{0.2308}}
  >{\raggedright\arraybackslash}p{(\columnwidth - 4\tabcolsep) * \real{0.5000}}@{}}
\toprule\noalign{}
\begin{minipage}[b]{\linewidth}\raggedright
Margin
\end{minipage} & \begin{minipage}[b]{\linewidth}\raggedright
Value
\end{minipage} & \begin{minipage}[b]{\linewidth}\raggedright
Rationale
\end{minipage} \\
\midrule\noalign{}
\endhead
\bottomrule\noalign{}
\endlastfoot
Consequential shift (Exps 1--2) & 3 percentage points &
Smallest change in a reported success rate treated as
consequential; the same margin in both experiments keeps one
scale \\
Large shift (P2, S6) & 5 points (support); 2 points
(absence) & 5 points is the smallest gap within a close pair
(5--10 points), so a shift of that size can reverse a close
pair when the two policies' shifts differ by at least their
gap, which S4 tests. Absence requires the shift to be below
2 points; the absence statistic is biased upward (§D), so
the call is conservative \\
Misordering (P3, S4) & 5 points of excess misordering &
About one flipped far pair in twenty (P3); S4 uses the same
5 points for close pairs \\
Recalibration (S3a, S3b) & ratio \ensuremath{<} 0.5 &
Recalibration must at least halve the level error to count
as a fix \\
Reuse gap (Exp 3) & 0.010 consequential; ±0.002 negligible &
On a test set with reuse fraction p, the error advantage of
the model continued on reused trajectories (arm C, §F) over
the one continued on independent events (arm B) is pG − H,
where H is the harm below. Reused test items raise this
advantage by pG, relative to an evaluation on independent
events only. At 10\% reuse, a 0.010 gap raises it by 0.001
arena widths, the ranking margin we set for closely matched
models (a stipulated value, not an estimate). At 50\% reuse,
a 0.002 gap does the same \\
Harm (Exp 3) & 0.001 & On independent events only, C's error
advantage over B is −H. Reuse must raise error on
independent events by less than 0.001 arena widths, the same
ranking margin \\
\end{longtable}
\endgroup

\needspace{6\baselineskip}
\hypertarget{a.-registered-code-freezes-and-amendments}{%
\section{A. Registered code, freezes and
amendments}\label{a.-registered-code-freezes-and-amendments}}

\begingroup\small
\begin{longtable}[]{@{}
  >{\raggedright\arraybackslash}p{(\columnwidth - 2\tabcolsep) * \real{0.2143}}
  >{\raggedright\arraybackslash}p{(\columnwidth - 2\tabcolsep) * \real{0.7857}}@{}}
\toprule\noalign{}
\begin{minipage}[b]{\linewidth}\raggedright
Item
\end{minipage} & \begin{minipage}[b]{\linewidth}\raggedright
Commitment
\end{minipage} \\
\midrule\noalign{}
\endhead
\bottomrule\noalign{}
\endlastfoot
Documents & The full protocol is
\texttt{docs/\allowbreak{}registration/\allowbreak{}protocol\_\allowbreak{}v3.4.md} at the
registered commit, cited below as ``Protocol''. It adds
operational detail, the evidence for each choice and the
outcome tables. It adds no confirmatory commitment. Where
the two differ, this document takes precedence and the
difference is logged as a deviation. The OSF Preregistration
(this document and the registration form) is embargoed until
peer review ends \\
Code at registration &
\texttt{physicalreasoning/\allowbreak{}trajectory-lineage} at tag
\texttt{osf-registration-v3.4}, the commit that adds this
document and the Protocol under \texttt{docs/\allowbreak{}registration/\allowbreak{}}
on top of \texttt{28e4ada}. The registration form records
the full hash of the tagged commit. Experiment 2 pipeline:
\texttt{subhashjprasad/\allowbreak{}robot-world-model-optimism-research}
at \texttt{\seqsplit{ed416e15ab9e4fd7c7b9ba4c51bdb343b4a28ce5}}. The
Experiment 2 and Experiment 3 analyzers are at the
registered commit (below). Experiment 1's analysis code is
recorded by the November 20 amendment, before any Experiment
1 outcome is computed. For every experiment, where code and
this document differ, this document governs \\
Code freezes & Each experiment receives a dated tag before
its confirmatory runs. Experiments 2 and 3 also have freeze
manifests that pin configuration, data selection, analyzer,
code and environment by hash. The paper lead and a
designated statistics reviewer (an author who runs no
confirmatory jobs) sign each manifest. The runners refuse an
unsigned manifest \\
November 20 amendment & Contains only these items.
\textbf{Experiment 1:} the pilot's operational values (token
grid, training steps, judge weights, which §D fallbacks
fired, gate results, ACT chunk size, the confirmatory
hardware class; the exploratory released world model, if
any; if the encoder arm runs, its readout and how its
2-frame tubelet steps align with per-step actions), the 12
policy hashes, the bootstrap percentiles from fallback 7
with the hash of the calibration script, and the repository,
commit and analysis code. \textbf{Experiment 3:} the
detector-audit details listed in §F. Filed before any
Experiment 1 outcome or 2D detector test outcome is
computed. The Experiment 1 analysis code checks the
amendment's hash before reading outcomes. The amendment
cannot change hypotheses, endpoints, thresholds, exclusions
or analysis families. If it is not filed before any
Experiment 1 outcome is computed, Experiment 1 becomes
exploratory; if it is not filed before any 2D detector test
outcome is scored, H3.3--H3.5 become exploratory \\
Analyzers at registration (Experiments 2 and 3) & The
registered commit contains the analyzers that implement §C's
Holm rule (Experiment 2) and §F (Experiment 3), with their
tests and rerun dry runs (code commit \texttt{f361b7e}). No
code amendment is needed for them. The runners refuse to
start without signed freeze manifests that pin those
analyzers \\
December 8 amendment & Execution report, deviations, and a
claims ledger with the exact abstract sentence for each
registered outcome. The headline sentence carries the P2
outcome and the qualifiers that §D requires. The amendment
is written before unblinding. It authorizes unblinding and
cannot select analyses after results \\
Amendment timing & Amendment dates are targets. The
Experiment 1 and Experiment 3 parts of the November 20
amendment may be filed separately; each is valid if filed
before the first outcome it governs is computed. The
December 8 amendment may be filed per experiment \\
Unblinding & The statistics reviewer unblinds each
experiment from its tag, only after that experiment's part
of the December 8 amendment is filed. Until then, run
operators are not shown outcome values: Experiment 1 pilot
reports show no world-model shift, and Experiment 2's runner
at its freeze tag logs no estimated success; its run files,
which hold per-episode judge scores, are analyzed only by
the statistics reviewer \\
\end{longtable}
\endgroup

\needspace{4\baselineskip}
\hypertarget{b.-hypothesis}{%
\section{B. Hypothesis}\label{b.-hypothesis}}

Estimated success from world-model policy evaluation depends
on the share of failures in the world model's training data,
while policies far apart in true success are misordered
little more often in imagined rollouts than by the same
judge on real rollouts. Reusing training trajectories in
evaluation lowers forecast error on those trajectories
relative to independent events. Experiment 1 (primary) tests
the failure-share hypothesis in closed-loop policy
evaluation. Experiment 2 checks the open-loop evidence on
logged robot data. Experiment 3 tests the reuse hypothesis
in controlled 2D forecasting; its effect on policy
evaluation is registered for a follow-up (§F).

\needspace{4\baselineskip}
\hypertarget{c.-analysis-rules-common-to-all-experiments}{%
\section{C. Analysis rules common to all
experiments}\label{c.-analysis-rules-common-to-all-experiments}}

\begin{itemize}
\tightlist
\item
  \textbf{Unit of replication.} One training run (seed) in
  Experiments 1--2; one block in Experiment 3. Policies,
  pairs, frames and renders are not replicates.
\item
  \textbf{Outcomes.} Each claim is \emph{supported} (the
  bound clears the threshold), \emph{absent} (the
  equivalence interval lies inside the margin),
  \emph{refuted} (the bound clears the threshold in the
  opposite direction), or otherwise \emph{inconclusive}.
  Missing a threshold never counts as absence. Decisions use
  bounds, not point estimates.
\item
  \textbf{Bounds.} Support uses one-sided 95\% bounds in
  Experiments 1--2 and the ends of two-sided 95\% intervals
  (one-sided 2.5\%) in Experiment 3, except its detector
  family (§F). Absence uses 90\% intervals (two one-sided
  tests at 5\% each).
\item
  \textbf{Intervals.} Paired t over seeds (Experiments 1--2)
  or blocks (Experiment 3); Welch t where Experiment 3
  compares separate groups of blocks; a bootstrap for
  Experiment 1's P2 and S6 (§D).
\item
  \textbf{Confirmatory families.} Experiment 1: headline (P1
  and P3), P2, and eight secondaries. Experiment 2: primary
  and ten secondary p-values. Experiment 3: primary claim,
  seven secondaries, and the detector family H3.3--H3.5 (if
  the November 20 amendment records their details before any
  2D detector test outcome is scored). Each family has its
  own error control, with no correction across families.
  Secondary families use Holm correction.
\item
  \textbf{Holm families.} In every Holm family, a member
  that is untested, descriptive, or inconclusive because of
  an unresolved run or too few screened blocks enters with p
  = 1, so each family keeps its registered size. A member is
  untested when its section says so (S5's checks, a J2 that
  fails even the 5 cm rule, a detector that fails the copy
  gate). Where a section defines a member's p-value (S6 in
  fallback 7; the detector family in §F), that definition
  applies; otherwise, when all seed or block values are
  equal, p is 0 if the common value lies beyond the
  threshold in the claimed direction and 1 otherwise. A
  member is declared when its Holm-adjusted p is at most the
  familywise level.
\item
  \textbf{Stopping rule.} Only the scheduled counts are run.
  There are no interim looks and no extension for
  significance. The hard budget cap is \usd{}2,000.
\item
  \textbf{Exclusion criteria.} Data are excluded only under
  frozen eligibility rules, for verified data corruption, or
  when evaluation items duplicate training items where the
  design does not intend it. Statistical outliers are
  retained. No exclusion is made after outcomes are visible.
\item
  \textbf{Missing runs.} Each experiment's section gives its
  retry and replacement rules and what an unresolved run
  implies. Where a section states nothing, each missing
  value is set to every possible value (rates from 0 to 1),
  and a decisive claim stands only if it holds for all of
  them; otherwise it is inconclusive.
\item
  \textbf{Non-confirmatory analyses.} Every analysis not
  listed as confirmatory is labelled exploratory or
  descriptive. An absence call means only ``no effect beyond
  the margin for that measure''; a difference between two
  measures (for example J1 vs J2) is read only from a paired
  contrast with its own interval, reported descriptively.
\item
  \textbf{Deviations.} Each departure from this document is
  logged with its date, reason, whether outcomes were
  accessible, and the claims it affects. A claim affected by
  a change made after its outcomes were accessible is
  exploratory, and the registered analysis is also reported
  where it can be run.
\item
  \textbf{Reporting.} All three experiments are reported
  regardless of outcome.
\end{itemize}

\needspace{4\baselineskip}
\hypertarget{d.-experiment-1-primary-closed-loop-evaluation-in-simulation}{%
\section{D. Experiment 1 (primary): closed-loop evaluation
in
simulation}\label{d.-experiment-1-primary-closed-loop-evaluation-in-simulation}}

\textbf{Setting.} ManiSkill3 PickCube-v1 with a Panda arm
and one fixed horizon. The horizon is the shortest of
50/\allowbreak 75/\allowbreak 100/\allowbreak 150/\allowbreak 200 steps at which a noiseless motion planner
succeeds in 95\% of 1,000 episodes. Success is read at the
final step. It requires the cube centre within 2.5 cm of the
goal and every arm joint moving at most 0.2 rad/s. The world
model is an action-conditioned latent transformer (width
256, depth 4). From the last 3 frames and the action, it
predicts the next frame's tokens and proprioception; it is
trained with a 4-step rollout loss. Per frame, it receives
17 frozen DINOv2-small tokens (the class token plus a 4×4
pooled patch grid), joint positions and velocities, and the
goal. An imagined rollout starts from one real frame, and
the policy then acts on the model's predictions until the
horizon.

\textbf{Intervention.} Experiment 1 varies the training
failure share φ ∈ \{0, .10, .25, .50, .65\}. Each arm draws
3,000 episodes, a fraction φ of them failures, from fixed
pools of 10,000 successes and 10,000 failures. Every
settings cell in the pools contributes equal numbers of
successes and failures. Cells below 10\% or above 90\%
success are dropped. A classifier that predicts the pool
from settings and start pose must have AUROC at most 0.55.
Otherwise, each failure is paired, without replacement, with
the nearest success in start pose within its settings cell.
Each arm has 12 seeds, in fixed halves 1--6 and 7--12. A
seed sets the initialization and a nested episode subsample.

\textbf{Policies.} 12 fixed policies, 4 from each family:
behaviour cloning, a PPO teacher distilled by DAgger, and
ACT (action-chunking transformer). Policies are selected on
true success only, before any world-model outcome. Selection
maximizes cross-family close pairs, subject to at least two
policies in each quarter of the success range and at least
20 far pairs. Ties go to the set with more far pairs. Pair
classes use each policy's 10,000-episode true success.

\textbf{Judges.} J1 (primary; the learned success judge) is
a small classifier with a fixed 0.5 threshold. It must reach
AUROC ≥ 0.80 and expected calibration error \ensuremath{<} 0.05
on real rollouts. J2 (check; the state-rule judge) is
trained without success labels. It decodes the cube position
and applies the success rule, using the world model's
predicted joint velocities in imagined rollouts. It must
reach median cube error \ensuremath{<} 1 cm and balanced
accuracy ≥ 0.90. If J1's AUROC is below 0.80, or its
expected calibration error is 0.05 or more after Platt
scaling (fallback 2), P1--P3 and the J1-based secondaries
are decided as registered but stated to apply only to a
judge that failed validation. \textbf{Ground truth:} 10,000
simulator episodes per policy. Each seed uses a fixed set of
2,000 starts.

Under the 5 cm rule (fallback 3), J2 applies the success
rule with a 5 cm distance threshold, while truth and J1 stay
at 2.5 cm. It passes if its balanced accuracy against true
success on real rollouts is at least 0.90; its median cube
error is reported and is not a pass criterion. Under this
rule S1 concerns success at 5 cm, so the share of episodes
whose decoded final cube distance lies between 2.5 and 5 cm
is reported for real and imagined rollouts at φ = .10 and
.65 (descriptive). J2's validation result fixes what S1 can
support:

\begingroup\small
\begin{longtable}[]{@{}
  >{\raggedright\arraybackslash}p{(\columnwidth - 4\tabcolsep) * \real{0.2250}}
  >{\raggedright\arraybackslash}p{(\columnwidth - 4\tabcolsep) * \real{0.3250}}
  >{\raggedright\arraybackslash}p{(\columnwidth - 4\tabcolsep) * \real{0.4500}}@{}}
\toprule\noalign{}
\begin{minipage}[b]{\linewidth}\raggedright
J2 validation
\end{minipage} & \begin{minipage}[b]{\linewidth}\raggedright
S1
\end{minipage} & \begin{minipage}[b]{\linewidth}\raggedright
Attribution of P1 to the world model
\end{minipage} \\
\midrule\noalign{}
\endhead
\bottomrule\noalign{}
\endlastfoot
Passes as registered (2.5 cm rule) & Tested as registered &
Made if S1 is supported (declared in the Holm family) \\
Passes only under the 5 cm rule & Tested, with its own
p-value in the Holm family; reported as ``S1 under the
relaxed 5 cm rule'' & Made if S1 is supported (declared in
the Holm family), and only with that qualifier in the same
sentence, in the abstract and the claims ledger. The
sentence states that the check used a 5 cm threshold where
the task uses 2.5 cm \\
Fails under the 5 cm rule & Untested; enters the Holm family
with p = 1 & Not made; stated as untested \\
\end{longtable}
\endgroup

\textbf{Estimands} (percentage points; per seed, policies
weighted equally). \emph{World-model shift:} the judge's
estimate on imagined rollouts minus its estimate on real
rollouts of the same policy from the same starts. The
difference removes the judge's bias on real rollouts; judge
error specific to imagined frames remains (see Limitations).
\emph{Excess misordering:} the rate at which imagined
estimates put a pair in the wrong order (ties count as
wrong) minus the same rate for the judge on real rollouts.
\emph{Absolute misordering:} the rate at which a set of
estimates puts a pair in the wrong order against true
success (ties count as wrong), computed for imagined
rollouts and for the judge on real rollouts. \emph{Far
pairs} differ by at least 15 points in true success;
\emph{close pairs} by 5 to 10. \emph{Margins:} 3 points is
the smallest consequential shift (as in Experiment 2). 5
points is the smallest close-pair gap, so a shift of that
size can reverse a close pair when the two policies' shifts
differ by at least their gap (S4).

\begingroup\small
\begin{longtable}[]{@{}
  >{\raggedright\arraybackslash}p{(\columnwidth - 6\tabcolsep) * \real{0.0667}}
  >{\raggedright\arraybackslash}p{(\columnwidth - 6\tabcolsep) * \real{0.2222}}
  >{\raggedright\arraybackslash}p{(\columnwidth - 6\tabcolsep) * \real{0.4000}}
  >{\raggedright\arraybackslash}p{(\columnwidth - 6\tabcolsep) * \real{0.3111}}@{}}
\toprule\noalign{}
\begin{minipage}[b]{\linewidth}\raggedright
ID
\end{minipage} & \begin{minipage}[b]{\linewidth}\raggedright
Claim
\end{minipage} & \begin{minipage}[b]{\linewidth}\raggedright
Supported if
\end{minipage} & \begin{minipage}[b]{\linewidth}\raggedright
Absent or refuted if
\end{minipage} \\
\midrule\noalign{}
\endhead
\bottomrule\noalign{}
\endlastfoot
P1 & Failure share moves estimates & Mean shift at φ = .10
minus φ = .65: lower bound \ensuremath{>} 3 & Absent: 90\%
interval inside ±3 \\
P2 & The world-model shift is large at a 10\% failure share
& Cross-fitted mean absolute shift at φ = .10: bootstrap
lower bound \ensuremath{>} 5 & Absent: bootstrap upper bound
\ensuremath{<} 2 for the mean over policies of each policy's
absolute 12-seed mean shift \\
P3 & Far pairs keep their order relative to the judge on
real rollouts & Far-pair excess misordering at φ = .10:
upper bound \ensuremath{<} 5 & Refuted: lower bound
\ensuremath{>} 5 \\
S1 & P1 holds with J2 & As P1, with J2 & As P1 \\
S2 & The shift falls with φ & Per-seed slope over φ = .10,
.25, .50: upper bound \ensuremath{<} 0 & n/a \\
S3a & One logit offset, fitted on four anchor policies,
fixes levels & Mean absolute error against truth on the
other 8 policies at φ = .10, after the fit over before:
upper bound \ensuremath{<} 0.5 & Refuted: lower bound
\ensuremath{>} 0.5 \\
S3b & An intercept and a slope fix levels & As S3a, with a
logistic intercept-and-slope refit (slope \ensuremath{>} 0)
& As S3a \\
S4 & Close comparisons are fragile & Close-pair excess
misordering at φ = .10: lower bound \ensuremath{>} 5 &
Absent: 90\% upper bound \ensuremath{<} 5 \\
S5 & P1 holds for a stochastic world model & As P1 (5
seeds). Counts only if the competence check passes and the
latent is used: on at least 5\% of held-out episodes, 20
samples from the same start and actions do not all get the
same J1 label & As P1 \\
S6 & The shift is large with no failures & As P2, at φ = 0 &
As P2 \\
S7 & P1 holds in a larger model trained on more data & Shift
at φ = .10 minus .65 for the width-512, 9,000-episode model
(5 seeds): lower bound \ensuremath{>} 3 & As P1 \\
\end{longtable}
\endgroup

\textbf{Decision rule.} \emph{Headline.} The headline holds
only if P1 and P3 are both supported, each at one-sided 5\%,
with no correction. \emph{P2.} P2 is decided separately at
its own one-sided 5\% and reported next to the headline; the
headline and P2 are separate claims, each controlling its
own error, and no joint rate is claimed. The headline does
not establish that the shift is large; only P2 does. The
abstract and the claims ledger state the P2 outcome in the
same sentence as the headline, in one of three fixed forms.
Supported: the mean absolute shift at a 10\% failure share
exceeds 5 points. Absent: it is below 2 points.
Inconclusive: its size is not established. \emph{P3
wording.} P3 bounds misordering in excess of the judge's on
real rollouts, so it can be supported when that judge itself
misorders far pairs. Absolute far-pair misordering against
true success at φ = .10, for imagined rollouts and for J1 on
real rollouts, is reported next to P3, each as the mean over
the 12 seeds with a two-sided 95\% t interval (descriptive).
Excess misordering is a difference of two rates, so the
share of far pairs that J1 orders correctly on real rollouts
and wrongly on imagined rollouts, and the reverse share, are
reported with it (descriptive). If J1's absolute far-pair
misordering on real rollouts is above zero in any of the 12
seeds, a supported P3 and a supported headline are worded as
limited additional ordering error in imagined rollouts. If
it is zero in every seed, they are worded as far-pair
misordering in imagined rollouts below 5 points against true
success. Neither is worded as correct ranking of far pairs.
\emph{Secondaries.} The secondary family is S1, S2, S3a,
S3b, S4, S5, S6 and S7. Holm correction applies to their
eight one-sided support p-values. A secondary that is
untested or descriptive (for example S5 when its usage check
fails, or S5 and S7 after fallback 6) enters with p = 1, so
the family always has eight members. S6's p-value is defined
in fallback 7. Secondary absence and refutation calls are
unadjusted, labelled as such, and enter the claims ledger
only as exploratory readings. \emph{Attribution and scope.}
A supported P1 shows that the estimate from the world model
and J1 together moves with the failure share. It is
attributed to the world model, and not to how J1 reads
imagined frames, only if S1 (the same test with J2) is also
supported. When J2 passed only under the 5 cm rule, the
attribution carries the qualifier in the Judges table. The
bounds describe these 12 policies and make no claim about a
population of policies.

\textbf{Estimators.} Each per-seed quantity has a t interval
over the 12 seeds (5 seeds for S5 and S7; S2: least-squares
slope; S3a and S3b: log ratio). If all seed values are
equal, the bound is that common value. P3 is also reported,
descriptively, with an exact sign-flip permutation bound
over seeds. P2 and S6 use split-seed cross-fitting. In each
seed, each policy's shift is signed by its mean shift in the
other half, and the signed shifts are averaged. The P2 and
S6 bounds come from 10,000 bootstrap draws that resample,
within each half, the seeds present in that half and
recompute every step, at the percentiles fixed by fallback
7. A policy whose other-half mean shift is exactly 0 gets
sign 0. For absence, P2 and S6 use the mean over policies of
each policy's absolute 12-seed mean shift, which is biased
only upward. The S3 anchors are the policies nearest the
20th, 40th, 60th and 80th percentiles of true success, each
with 100 labelled real rollouts drawn, for each seed, from
that seed's own real rollouts. Both fits use logit(real
success) = a + b·logit(imagined success), by binomial
maximum likelihood over each anchor's 100 labels: S3a fits a
with b = 1; S3b fits a and b, with b \ensuremath{>} 0.
Imagined success rates are clipped to {[}1/4,000,
3,999/4,000{]} before any logit is taken.

\textbf{Permitted fallbacks} (each recorded in the November
20 amendment):

\begin{enumerate}
\def\labelenumi{\arabic{enumi}.}
\tightlist
\item
  If the pilot competence check fails, use the 16-step
  rollout loss first, then the 8×8 token grid (65 tokens).
  If the behaviour-cloning check fails, use the 8×8 grid. A
  grid change requires retraining and revalidating both
  judges and rerunning the pilot.
\item
  If J1's expected calibration error is 0.05 or more, apply
  Platt scaling.
\item
  If J2 fails validation, use the 8×8 grid, then the 5 cm
  rule for J2 only. The table under Judges fixes what S1 can
  then support.
\item
  Policy selection: first lower the far-pair minimum from 20
  to 15, then allow a single policy in one success quarter.
  If the third family is not ready when policies are
  selected, use 8 policies from two families (two per
  quarter, at least 10 far pairs), run S3a with 2 anchors
  (33rd and 67th percentiles), and make S3b descriptive.
\item
  If a competence gate still fails at the last setting,
  P1--P3 are decided as registered but stated to apply only
  to a world model that failed that gate.
\item
  If Experiment 1 would exceed its 1,000 GPU-hour cap, drop
  the optional PushCube-v1 replication, then the V-JEPA 2
  encoder arm (descriptive; below), then cut S7 and S5 to 3
  seeds, which makes them descriptive. If the main sweep
  alone would exceed the cap, the design is amended as a
  deviation before it runs.
\item
  The pilot's synthetic simulation sets the bootstrap
  percentiles before the main sweep. It is run once, with
  20,000 replicates per scenario, 10,000 fresh bootstrap
  draws in each replicate, and a random seed written in the
  script before the run.

  \begin{itemize}
  \tightlist
  \item
    \emph{(a) Scenarios.} Every scenario holds the mean
    absolute true shift at the margin (5 for support, 2 for
    absence). Five patterns set the policies' absolute
    shifts, as multiples of the margin: equal (every policy
    at 1); a quarter at zero (a quarter of the policies at
    0, the rest at 4/3); graded (policy \(k\) of \(n\) at
    \(2(k-1)/(n-1)\)); concentrated (a quarter at 2.8, the
    rest at 0.4); and half near zero (half at 0.1, half at
    1.9). Each pattern is run twice: with signs alternating
    across policies ordered by absolute shift, and with
    every sign positive. The scenarios also vary the
    per-policy, per-seed SD of the shift (1 and 3), the
    cross-policy correlation ρ (0 and 0.5) and the error
    distribution (normal and t(3)). Each error is
    \(\mathrm{SD} \times (\sqrt{\rho}\,c + \sqrt{1-\rho}\,e)\),
    independent across seeds, where \(c\) is shared by the
    policies in a seed, \(e\) is specific to a policy, and
    both are standard normal or t(3) divided by
    \(\sqrt{3}\). That gives 80 scenarios, each run at both
    margins; the same patterns apply to 8 policies under
    fallback 4.
  \item
    \emph{(b) Pass rule.} Coverage is the share of
    replicates in which the support lower bound is at or
    below the margin, or the absence upper bound is at or
    above it. Percentiles are searched on a 0.25-point grid.
    To allow for Monte Carlo error, a percentile passes a
    scenario only if the lower end of the one-sided 95\%
    Clopper--Pearson interval for its coverage in that
    scenario is at least 95\%.
  \item
    \emph{(c) Percentile for the bounds.} The decision uses
    the least extreme percentile that passes every scenario;
    it is never less extreme than the 5th (support) or 95th
    (absence) and is capped at the 1st or 99th. A bound that
    does not pass at the cap cannot establish support or
    absence.
  \item
    \emph{(d) S6's Holm p-value.} The same simulation
    calibrates, for each Holm level \(\alpha_j = 0.05/j\) (j
    = 1, \ldots, 8), the least extreme percentile on the
    grid that is no less extreme than \(100\cdot\alpha_j\),
    is no further out than the 0.25th, and passes every
    scenario with required coverage \(1 - \alpha_j\) (the
    lower end of the same one-sided 95\% Clopper--Pearson
    interval is at least \(1 - \alpha_j\)). Calibrated
    percentiles are monotone in \(\alpha_j\); a level with
    no such percentile is unreachable. S6's p-value is the
    smallest reachable \(\alpha_j\) at which its lower bound
    exceeds 5, or 1 if there is none. In a planning run the
    0.25th percentile covered about 0.978 in the worst
    scenario, so S6 is expected to be declarable only late
    in the Holm sequence (Protocol §1.5).
  \end{itemize}
\end{enumerate}

\textbf{Competence gates.} The pilot gates use open-loop
replays averaged over 3 seeds. The decoded final cube
distance must separate real successes from failures with
AUROC ≥ 0.75. Predicted final stillness must reach balanced
accuracy ≥ 0.85. A quick behaviour-cloning run must reach
75\% success. Each world-model gate must pass at φ = .10 and
at φ = .65 separately. S5's competence check applies the
same gates to the stochastic model at φ = .10 and .65,
averaged over its first 3 seeds.

\textbf{Failed runs.} An infrastructure failure
(input/output, memory, device or preemption error) at any
stage is retried once with the identical seed: training
resumes from the last complete checkpoint, and rollouts and
scoring are rerun from the trained checkpoint on the same
starts. A training failure (a non-finite loss, or a crash
with no infrastructure cause) that is detected before any
rollout of that run is scored reruns its seed index once, in
every failure share of its model configuration (five for the
main sweep and S5, two for S7), under the next unused
reserve index (13, 14, 15 for seeds 1--6; 16, 17, 18 for
seeds 7--12). The rerun keeps its half and that half's
evaluation starts, and gets the same infrastructure retry.
No other failure is replaced. Any run that remains
unresolved, including a failed retry, a failed reserve run
or a training failure found after scoring has started,
leaves its index missing, and every decision that uses it is
inconclusive. A sensitivity analysis drops every replaced
index from all arms. If the primary analysis supports a
claim and the sensitivity analysis does not, the claim is
reported as ``supported (not robust to replacement)''. It
counts as supported in the decision rules, but neither it
nor a headline that relies on it is stated as supported in
the abstract. Any other disagreement is reported beside the
primary call, which stands. Invalid imagined rollouts are
scored as failed episodes.

\textbf{Descriptive encoder arm} (run only if the pilot's
cost projection allows). The main world model at φ = .10 and
.65 (5 seeds) runs on a frozen V-JEPA 2-L encoder. A
readout, fitted once on the full balanced pools and not per
arm, maps the encoder output to the 17 DINOv2-small tokens
that policies and judges read. The arm's reference is real
rollouts in which policies and judges read the readout's
tokens of real frames, so the readout's error on real frames
cancels in the shift, as the judge's bias does. This arm's
shift is reported next to the main model's. DINOv2-small
stays the main encoder because Experiment 2 uses it.

\textbf{Replay check} (descriptive). For φ = .10 and .65,
the action sequence of each imagined rollout on the first
200 starts of each seed is executed open loop in the
simulator from the same start state. We report the true
success of these executions, the error of the imagined final
cube position against the executed one, and J1 on the
executed frames. The check is reported only if replaying the
recorded actions of real rollouts reproduces their outcomes
in at least 99\% of episodes.

\textbf{Limitations stated with every result.} World-model
error and judge error are separated only in part: the shift
removes the judge's bias on real rollouts, and the replay
check measures the world model's error on the action
sequences policies produced in imagination, but judge error
specific to imagined frames remains. Experiment 1 uses one
small world model on frozen DINOv2-small tokens, one
simulated task and 12 fixed policies; support would not by
itself extend to other tasks, to larger or pixel-generative
video world models, or to planning with a world model.
``Failure share'' means the resampling intervention above,
which also changes coverage. The study uses one fixed
horizon, set by the rule in Setting (expected to be about
100 steps). The distilled-PPO family imitates the teacher
that generated half the training data, so it is closer to
that data than the other families are. Results are also
reported per family.

\needspace{4\baselineskip}
\hypertarget{e.-experiment-2-transition-matched-kuka-control}{%
\section{E. Experiment 2: transition-matched Kuka
control}\label{e.-experiment-2-transition-matched-kuka-control}}

\textbf{Data.} Logged Kuka episodes from Open X-Embodiment,
rebuilt from the public shards with the original pipeline's
extractor (episode manifest \texttt{1dcb390f…}). This
rebuilt copy is the registered data. Prior results on these
data are known (§G). The world model and features are the
original pipeline's, unchanged, trained for 2,000 steps per
run. The success judge is fitted once and frozen. If its
AUROC is below 0.75, the runs do not start and Experiment 2
reports no confirmatory result.

\textbf{Designs.} Runs use seeds 100--109, paired across
arms. Historical seeds 0--2 are excluded. Every selection is
certified before training.

\begin{itemize}
\tightlist
\item
  \textbf{O:} the original 3,000-episode arms at φ = .10 and
  .65 (20 runs).
\item
  \textbf{A (primary):} exactly 31,600 distinct training
  transitions at φ ∈ \{0, .10, .25, .50, .65\} (50 runs).
\item
  \textbf{B:} φ = .10 and .65, with the same normalized
  length distribution for successes and failures, at 31,600
  transitions (20 runs).
\item
  \textbf{V:} a shared 1,500-success base plus further
  successes (V+S) or failures (V+F) up to exactly 31,600
  transitions (20 runs).
\end{itemize}

Total: 110 runs, or 130 under the B fallback.

\textbf{Estimand.} For design d,
\(\Delta_d = 100 \times (R_{d,.10} - R_{d,.65})\), averaged
over seeds. R is the estimated success rate on evaluation
episodes, from the judge's thresholded prediction (1 if p
\ensuremath{>} 0.5) on each imagined episode. Episodes are
weighted so that each retained length's total weight is
proportional to the smaller of its shares among retained
successes and among retained failures, and within each
length the successes carry 0.355797 of that weight (the
success share among all 2,760 eligible evaluation episodes)
and the failures the rest. Length therefore carries no
outcome information. For V,
\(\Delta_V = 100 \times (R_{V+S} - R_{V+F})\).

\begingroup\small
\begin{longtable}[]{@{}
  >{\raggedright\arraybackslash}p{(\columnwidth - 4\tabcolsep) * \real{0.2571}}
  >{\raggedright\arraybackslash}p{(\columnwidth - 4\tabcolsep) * \real{0.3714}}
  >{\raggedright\arraybackslash}p{(\columnwidth - 4\tabcolsep) * \real{0.3714}}@{}}
\toprule\noalign{}
\begin{minipage}[b]{\linewidth}\raggedright
Quantity
\end{minipage} & \begin{minipage}[b]{\linewidth}\raggedright
Supported if
\end{minipage} & \begin{minipage}[b]{\linewidth}\raggedright
Absent if
\end{minipage} \\
\midrule\noalign{}
\endhead
\bottomrule\noalign{}
\endlastfoot
\(\Delta_A\) (primary) & One-sided 95\% lower bound
\ensuremath{>} 3 & 90\% interval inside (−3, 3) \\
\(\Delta_O\), \(\Delta_B\), \(\Delta_V\) & Holm-adjusted
support test of a drop \ensuremath{>} 3 & Holm-adjusted
equivalence test within ±3 \\
\(\Delta_O\) − \(\Delta_A\), \(\Delta_A\) − \(\Delta_B\) &
Holm-adjusted support test of an attenuation \ensuremath{>}
3 & Holm-adjusted equivalence test within ±3 \\
\end{longtable}
\endgroup

\textbf{Decision rule.} Tests use paired t over seeds, with
two one-sided tests for equivalence. \(\Delta_A\) is the
primary quantity, tested at one-sided 5\%. The secondary
family has ten p-values (support and equivalence for each of
the five secondary quantities), with Holm correction at
familywise 0.05. A secondary is declared when its
Holm-adjusted p is at most 0.05. Each interval is reported
at level \(1 - 2\alpha_j\), where \(\alpha_j\) is that
hypothesis's Holm step level.

\textbf{Eligibility and exclusions.} Only episodes of 3 to
24 frames are eligible, within the existing shard-based
evaluation split. A length enters the evaluation only if
both training outcomes occur at that length and it has at
least 20 evaluation episodes per outcome. The set of
eligible lengths is fixed from labels and lengths before any
new score. Difficult episodes and poor seeds are never
excluded.

\textbf{Permitted fallback.} If 31,600 transitions is
infeasible for B, B uses the largest multiple of 1,000 not
above 0.9 × the maximum feasible budget, and A comparators
are added at that budget (130 runs).

\textbf{Failed runs.} Infrastructure failures are retried
once with the same seed. An unresolved run makes every
decision that uses it inconclusive.

\textbf{Descriptive analyses.} Mean-probability response,
intermediate failure shares, Brier scores and AUROC.

\needspace{4\baselineskip}
\hypertarget{f.-experiment-3-trajectory-reuse-and-duplicate-detector-audit}{%
\section{F. Experiment 3: trajectory reuse and
duplicate-detector
audit}\label{f.-experiment-3-trajectory-reuse-and-duplicate-detector-audit}}

\textbf{Setting.} Two equal disks in a unit arena collide
elastically once; the contact always falls inside the
16-frame forecast window. A predictor sees 8 frames and
forecasts both disks' positions for the next 16. Error is
the mean over future disk-frames of min(1, distance), in
arena widths. Non-finite outputs score 1 and are kept.

\textbf{Block.} A base predictor trained on 512 events, then
five matched continuations of it:

\begin{itemize}
\tightlist
\item
  B: independent events (control);
\item
  C: the 64 evaluation-source trajectories with new
  appearances (reuse);
\item
  D: independent neighbors closer than any base-training
  event;
\item
  D2: a distance reference (descriptive);
\item
  E: the exact evaluation renders (positive control).
\end{itemize}

\textbf{Screening.} A block enters the primary analysis only
if its base model beats constant-velocity extrapolation on
256 separate screen events. Failed screens are reported,
with an all-block sensitivity analysis.

\textbf{Estimands.} L(X, y) is the error of the model
continued in arm X on test set y, where s = unseen new
renders of the source trajectories, i = independent events,
x = exact renders. Source gain S = L(B,s) − L(C,s); harm H =
L(C,i) − L(B,i); \textbf{reuse gap G = S + H}; event
specificity Q = {[}L(D,s) − L(C,s){]} − {[}L(D,i) −
L(C,i){]}; exact gain E = L(B,x) − L(E,x). The quantities
lower95 and upper95 are the ends of the two-sided 95\%
interval over blocks.

\textbf{Sample sizes.} The main confirmation has 1,600
ConvLSTM blocks: 800 at 16 hidden channels (primary), 400 at
64 and 400 at 256. Four axis configurations have 200 blocks
each: a transformer that regresses future disk positions
from frozen DINOv2-small or V-JEPA 2-L tokens (it predicts
coordinates, not future latent states), and 16 or 256 reused
sources. A decisive claim requires at least 720 screened
blocks at the primary level, 360 at each larger level and
180 per axis configuration.

\textbf{Margins} (arena widths): consequential 0.010;
negligible ±0.002; harm 0.001. \emph{Prevalence} p is the
share of evaluation items that are reused. On a test set
with prevalence p, C's error advantage over B is pG − H. The
term pG is the increase in this advantage relative to an
evaluation on independent events only. It is not C's full
advantage. At 10\% reuse, a 0.010 gap gives an increase of
0.001, the ranking margin we set for closely matched models.
At 50\% reuse, a 0.002 gap gives the same increase. The
order of C and B reverses between the two evaluations only
if H lies strictly between 0 and pG (0 \ensuremath{<} H
\ensuremath{<} pG when G is positive). In the claim labels
below, ``reuse advantage'' names the gap G, not C's full
advantage pG − H.

\textbf{Primary claim} (16-channel blocks only):

\begingroup\small
\begin{longtable}[]{@{}
  >{\raggedright\arraybackslash}p{(\columnwidth - 2\tabcolsep) * \real{0.4000}}
  >{\raggedright\arraybackslash}p{(\columnwidth - 2\tabcolsep) * \real{0.6000}}@{}}
\toprule\noalign{}
\begin{minipage}[b]{\linewidth}\raggedright
Claim
\end{minipage} & \begin{minipage}[b]{\linewidth}\raggedright
Requires
\end{minipage} \\
\midrule\noalign{}
\endhead
\bottomrule\noalign{}
\endlastfoot
Consequential reuse advantage & lower95(G) \ensuremath{>}
0.010, lower95(S) \ensuremath{>} 0 and upper95(H)
\ensuremath{<} 0.001 \\
Prevalence-dependent reuse advantage & 0.002 \ensuremath{<}
lower95(G) and upper95(G) ≤ 0.010, lower95(S) \ensuremath{>}
0 and upper95(H) \ensuremath{<} 0.001 \\
Detectable, magnitude unresolved & 0.002 \ensuremath{<}
lower95(G) ≤ 0.010 \ensuremath{<} upper95(G), lower95(S)
\ensuremath{>} 0 and upper95(H) \ensuremath{<} 0.001 \\
Gap detected, safeguards not established & lower95(G)
\ensuremath{>} 0.002 with lower95(S) ≤ 0 or upper95(H) ≥
0.001 \\
Negligible differential gap & 90\% interval of G inside
(−0.002, 0.002) \textbf{and} lower95(E) \ensuremath{>}
0.002 \\
Negligible gap, control below the margin & Negligible G and
0 \ensuremath{<} lower95(E) ≤ 0.002: inconclusive about
negligibility \\
Positive control not detected & Negligible G and lower95(E)
≤ 0 \\
Reverse gap & upper95(G) \ensuremath{<} −0.002 \\
Inconclusive & Otherwise, or any unresolved
main-confirmation block failure, or fewer screened blocks
than the minimum \\
\end{longtable}
\endgroup

\textbf{Secondary hypotheses.} One Holm family at familywise
one-sided 2.5\%, always seven p-values:

\begin{enumerate}
\def\labelenumi{\arabic{enumi}.}
\tightlist
\item
  Event-specific: Q \ensuremath{>} 0.010.
\item
  Detectable: G \ensuremath{>} 0.
\item
  Consequential harm: H \ensuremath{>} 0.001.
\item
  Capacity increase: G at 256 channels minus G at 16
  \ensuremath{>} 0 (Welch).
\item
  Reuse detectable with the frozen DINOv2 encoder: G
  \ensuremath{>} 0 in that configuration.
\item
  Reuse detectable with the frozen V-JEPA 2 encoder: G
  \ensuremath{>} 0 in that configuration.
\item
  Dose increase: G with 16 sources minus G with 256
  \ensuremath{>} 0 (Welch).
\end{enumerate}

Secondaries 1--3 use the screened 16-channel blocks;
secondary 4 compares the screened 256- and 16-channel
blocks; secondaries 5--7 use the screened blocks of their
configurations. The primary claim never depends on an axis.
An axis configuration with an unresolved block failure or
too few screened blocks enters with p = 1 and is labelled
unresolved.

\textbf{Labels.} A difference is labelled \emph{capacity-}
or \emph{dose-independent} if its 90\% interval lies inside
±0.002 and lower95(E) \ensuremath{>} 0.002 in each compared
configuration. A configuration is labelled \emph{negligible
with that encoder} if its 90\% interval of G lies inside
±0.002 and its lower95(E) \ensuremath{>} 0.002. Otherwise
the label is unresolved.

\textbf{Exploratory encoder comparison.} G with V-JEPA 2
minus G with DINOv2 (Welch interval), outside the Holm
family. It is reported as an encoder difference. It is not
attributed to pretraining alone, because the two encoders
also differ in size and in image vs video input.

\textbf{Detector audit.} The calibration and test splits
each have a 500-event gallery, 200 source trajectories under
8 transformations, and 600 ordinary and 200 close
independent events. The detectors are exact hash, SSCD (a
copy-detection embedding), DINOv2, OFS-3 (optical-flow
similarity), object-track matching, and true-state distance
(reference only). \emph{False removal} is the share of
independent events wrongly flagged. A detector counts only
if it first catches ordinary re-encoded copies (AUROC ≥ 0.95
against ordinary independent events and recall ≥ 0.80 at 5\%
false removal). A detector that fails is reported, and every
prediction involving it enters the family with p = 1.
Thresholds at 1\% and 5\% false removal are set on the
calibration split and frozen before test scoring. Intervals
and p-values come from a 10,000-draw bootstrap over
calibration and test source families; in each draw the
thresholds are re-estimated on the resampled calibration
negatives.

\emph{Appearance-only rerenders} are the new-appearance
transformation (transformation 2 of the 8). The copy gate is
judged on the calibration split. Each prediction is tested
with a one-sided bootstrap p-value, p = (1 + number of draws
on the null side of the threshold) / 10,001, where a draw
equal to the threshold counts as the null side; a compound
prediction takes the largest of its parts' p-values.

\begin{itemize}
\tightlist
\item
  H3.3: SSCD and DINOv2 each catch fewer than 50\% of
  appearance-only rerenders at 1\% false removal.
\item
  H3.4: object-track matching catches more than 80\% of them
  at 1\% false removal, and more than OFS-3 does.
\item
  H3.5: at 1\% false removal on ordinary events,
  object-track matching removes close independent events at
  a higher rate than ordinary ones.
\end{itemize}

H3.3--H3.5 form one Holm family at familywise one-sided 0.05
(not the 2.5\% of §C). They are confirmatory only if the
November 20 amendment records exactly these details before
any 2D detector test outcome is scored: the threshold
quantiles and tie rule, the resampling of negative families,
and extraction-failure handling (a failure counts as no
match and stays in the denominator). Otherwise they are
exploratory.

\textbf{Filtering} (descriptive; no pass or fail).
Trajectory-level filtering at the frozen 5\% false-removal
threshold is reported only if it keeps at least 90\% of
independent items. It is compared against random removal of
the same size. The reuse gap among retained items is
reported against coverage.

\textbf{Failed runs.} Infrastructure failures are rerun once
with identical seeds. Algorithmic failures are outcomes and
are never replaced. Any unresolved main-confirmation block
failure makes the primary claim and every secondary that
uses main blocks inconclusive. Estimates are still reported.
An axis failure affects only that configuration (above).
Neither type of failure affects the detector audit.

\textbf{Limitations stated with every result.} The
Experiment 3 predictors forecast positions from frames. They
take no actions, are not rolled out on their own
predictions, and score no policy. Reuse enters through a
100-update continuation, not pretraining. Results concern
forecast error on synthetic single-collision events; they
are not evidence about physical reasoning in general or
about robot world-model evaluation.

\textbf{Sealed 1D confirmation.} The 352-block 1D
confirmation (run October 20) stays sealed until the 2D
reuse outcomes are unsealed. It is then analyzed once with
the analyzer frozen at \texttt{a25d881} and reported
separately, regardless of outcome. Its protocol was frozen
by the Experiment 3 lead alone and was not preregistered.

\textbf{Registered for a follow-up} (not a claim of this
study). In the Experiment 1 setup, reuse counts as
consequential if it shifts estimated success by at least 3
points with far-pair order kept, or if reuse that favours
one policy reorders a pair at least 5 points apart.

\needspace{4\baselineskip}
\hypertarget{g.-prior-data-and-exposure}{%
\section{G. Prior data and
exposure}\label{g.-prior-data-and-exposure}}

\begin{itemize}
\tightlist
\item
  \textbf{Experiment 1:} prospective. The pilot has not been
  run.
\item
  \textbf{Experiment 2:} the existing data were analyzed
  before registration. Known results: about 51\% estimated
  success (zero-failure arm) vs 36\% recorded; length alone
  predicts outcome about as well as the world model (AUROC
  0.764 vs 0.761); and a 10.3-point drop from 10\% to 65\%
  training failures. The new runs analyze previously
  examined data prospectively and are not a replication on
  untouched data. The registered data are the copy rebuilt
  from the public shards (2,760 evaluation episodes, 35.58\%
  success). This copy may differ slightly from the cache
  behind the original result, so comparisons with the
  historical 10.3-point drop are approximate. Experiment 2
  is therefore a robustness check of a known effect; on its
  own it cannot confirm the hypothesis.
\item
  \textbf{Experiment 3:} 1D pilots exposed dispersions and a
  coarse positive-control flag only. Lovejesh Bangar, who
  leads Experiment 3, has seen the 16-block 1D interim
  effect estimates (October 19), including an exploratory
  filtering analysis on 8 of those blocks. The 1D detector
  benchmarks and the exploratory 1D filtering analysis
  (which together informed H3.3--H3.5 and the descriptive
  filtering analysis) and the memory checks with a lookup
  predictor that memorizes training trajectories have also
  been seen. The 1D confirmation is sealed. Every 2D pilot
  release (October 23, October 27--28 capacity, October 28
  architecture and dose) contained dispersions, counts,
  runtime and one coarse positive-control pass/fail flag.
\item
  \textbf{Experiment 3 repository record:} the 16-block 1D
  interim estimates and the 8-block filtering analysis were
  published in the repository on October 19: the estimates
  and per-block values in
  \texttt{results/\allowbreak{}exposure\_\allowbreak{}interim\_\allowbreak{}analysis.json} and
  \texttt{results/\allowbreak{}exposure\_\allowbreak{}interim\_\allowbreak{}raw.json}
  (\texttt{af2c376}); the estimates, including the positive
  control, and the filtering analysis in FINDINGS §7 and
  \texttt{results/\allowbreak{}exposure\_\allowbreak{}interim\_\allowbreak{}filtering.json}
  (\texttt{9fa1778}); and a README summary of both
  (\texttt{5461ca6}). These interim results include the
  positive-control estimate. The 1D detector benchmarks and
  the lookup-predictor memory checks are also public. The 2D
  pilot reports in the repository contain only dispersions,
  counts, runtime and the coarse flag.
\item
  \textbf{Positive-control margin.} The 0.002 margin that
  the positive control must clear (§F) is the negligible
  bound fixed on October 16 (\texttt{f0e30eb}), before any
  1D effect estimate. The requirement that the control clear
  it comes from the October 23 design plan.
\end{itemize}

\textbf{Experiment 3 attestation.} Lovejesh Bangar, who
leads Experiment 3, attests that:

\begin{enumerate}
\def\labelenumi{\arabic{enumi}.}
\tightlist
\item
  the only learned-model reuse-effect values he has seen, or
  has shared with any author, are the 1D interim estimates
  above;
\item
  the 352-block 1D confirmation has never been analyzed;
\item
  the 2D rule changes follow from design logic and not from
  any 1D value;
\item
  no Experiment 3 effect value has been sent to an external
  reviewer. The 1D interim estimates are public, so any
  reader of the repository may have seen them.
\end{enumerate}

Anyone exposed to an interim summary is treated as informed.
Protocol §3.8 gives the timeline with dates and commits.

\textbf{Provenance.} Commits cited above: \texttt{a25d881}
(1D analyzer freeze, October 19); \texttt{f0e30eb}
(negligible bound, October 16); \texttt{9fa1778},
\texttt{5461ca6}, \texttt{af2c376} (1D interim publication,
October 19); \texttt{28e4ada} (repository base for the
registered commit); \texttt{f361b7e} (Experiment 2 and 3
analyzers);
\texttt{\seqsplit{ed416e15ab9e4fd7c7b9ba4c51bdb343b4a28ce5}}
(Experiment 2 pipeline). Experiment 2 episode manifest:
SHA-256
\texttt{\seqsplit{1dcb390f0f095205fb4b8d68d023b5dcdc94f378564423059e8be07dd54c09d8}}.

\end{document}
